(1) 若入射二次谐波为零，$A_2(0)=0$，对$z$积分得$A_2(L)=i\kappa A_1^2(e^{i\Delta kL}-1)/(i\Delta k)=i\kappa A_1^2 L e^{i\Delta k L/2}\operatorname{sinc}(\Delta k L/2)$，其中$\operatorname{sinc}x=\sin x/x$。在$\Delta k=0$处须取连续极限$A_2=i\kappa A_1^2 L$，而不是分母为零后宣布发散。
(2) 在相同模式规范下生成强度$I_2(L)\propto|A_2|^2=\kappa^2|A_1|^4L^2\operatorname{sinc}^2(\Delta kL/2)$。相位差积累到$\pi$的距离$L_c=\pi/|\Delta k|$称相干长度，超出后不同晶体段的新增二次谐波会部分相消。对精确相位匹配$\Delta k=0$，强度按$L^2$增强。
(3) 令非线性系数按符号函数$s(z)=\operatorname{sgn}[\cos(\Delta k z)]$周期翻转，其傅里叶展开在与失配相位匹配的基频处的复系数模为$2/\pi$。在跨越多个周期且忽略端部修正时，积分$s(z)e^{i\Delta kz}$的共振平均约$(2/\pi)L$，故有效系数$\kappa_{\rm eff}=(2/\pi)\kappa$，强度恢复近似$L^2$增长但相对完美天然相位匹配效率减少因子$4/\pi^2$。
(4) 如果二次谐波已经明显抽取泵浦能量，$A_1$不再常量，需联立$\chi^{(2)}$三波耦合方程；一个二次谐波光子能量$2\hbar\omega$，所以理想频率转换对应消耗两个基频光子，Manley–Rowe不变量需按光子流而非只比较电场幅值。有限群速度失配、衰减及横向衍射还会使长晶体的$L^2$增长失效。
